\documentclass[11pt,a4paper]{report}

\usepackage[utf8]{inputenc}
\usepackage[T1]{fontenc}
\usepackage[english]{babel}
\usepackage{lmodern}
\usepackage{geometry}
\usepackage{microtype}
\usepackage{array}
\usepackage{booktabs}
\usepackage{longtable}
\usepackage{tabularx}
\usepackage{enumitem}
\usepackage{hyperref}
\usepackage{xcolor}
\usepackage{listings}

\geometry{margin=1in}
\setlist{nosep,leftmargin=*}
\setlength{\parskip}{0.45em}
\setlength{\parindent}{0pt}
\emergencystretch=2em

\hypersetup{
    colorlinks=true,
    linkcolor=blue!50!black,
    urlcolor=blue!50!black,
    pdftitle={Methodological Support Report for the SDR Spectrum Monitoring Sensor Node},
    pdfauthor={OpenAI Codex}
}

\newcolumntype{Y}{>{\raggedright\arraybackslash}X}

\lstset{
    basicstyle=\ttfamily\small,
    breaklines=true,
    columns=fullflexible,
    frame=single,
    framerule=0.2pt
}

\newcommand{\module}[1]{\path{#1}}
\newcommand{\approach}[1]{\paragraph{Technical approach.} #1}
\newcommand{\scopep}[1]{\paragraph{Scope.} #1}
\newcommand{\realization}[1]{\paragraph{Code realization.} #1}

\begin{document}

\hypersetup{pageanchor=false}
\begin{titlepage}
\centering
{\Large Methodological Support Report\par}
\vspace{0.7cm}
{\huge SDR Spectrum Monitoring Sensor Node\par}
\vspace{1.0cm}
{\large Repository-grounded architectural and design analysis\par}
\vfill
\begin{tabular}{ll}
Artifact type: & Technical methodological report \\
System class: & Embedded cyber-physical SDR sensing node \\
Primary implementation languages: & C and Python \\
Compilation baseline: & CMake + Python virtual environment \\
Document date: & \today \\
\end{tabular}
\vfill
\begin{minipage}{0.88\textwidth}
\small
This document was prepared from direct inspection of the repository source code, build scripts,
runtime entry points, documentation stubs, and verification assets. The report therefore describes
the implemented system as it exists in the inspected codebase, rather than an idealized target
architecture.
\end{minipage}
\end{titlepage}
\hypersetup{pageanchor=true}

\pagenumbering{roman}

\chapter*{Abstract}
\addcontentsline{toc}{chapter}{Abstract}
This report provides methodological support for the SDR Spectrum Monitoring Sensor node implemented
in the repository. The analysis frames the node as a layered measurement system that combines
deterministic radio-frequency acquisition in C with supervisory orchestration in Python. The report
documents the architectural foundation, allocation of responsibilities, inter-process integration
mechanisms, resilience patterns, deployment model, and current verification evidence. Each analytical
section explicitly states its scope, technical basis, and module-level realization so that the text
can serve both as formal design documentation and as a starting point for future maintenance,
auditing, and extension activities.

\tableofcontents

\pagenumbering{arabic}

\chapter{Purpose and Methodological Basis}

\section{Document Objective}
\scopep{The purpose of this report is to formalize the engineering rationale behind the sensor node,
with emphasis on architectural decomposition, module responsibilities, design trade-offs, and
evidence of implementation. The report is not limited to user-facing behavior; it addresses the
full technical chain from hardware control and signal processing to orchestration, telemetry,
resilience, and deployment.}

\approach{The document follows a systems-engineering method based on three principles:
\begin{enumerate}
\item requirement-oriented decomposition of the sensing problem into acquisition, processing,
coordination, telemetry, and recoverability concerns;
\item architectural allocation of each concern to the most appropriate runtime domain according to
timing sensitivity, hardware coupling, and operational risk;
\item traceability from design intent to concrete source modules and scripts.
\end{enumerate}
This method is appropriate for mixed embedded systems because it avoids both extremes of
over-centralization and undocumented ad hoc growth.}

\realization{The repository already exhibits a strong separation of concerns that supports this
method: real-time RF execution is concentrated in \module{rf/rf.c} and \module{rf/libs/*},
telecommunication and positioning are handled in \module{gps-lte/gps-lte.c} and
\module{gps-lte/libs/*}, orchestration is centered in \module{orchestrator.py} and
\module{functions.py}, and cross-cutting infrastructure is concentrated in \module{cfg.py},
\module{utils/io\_util.py}, \module{utils/request\_util.py}, and \module{utils/status\_util.py}.}

\section{Analytical Evidence Base}
\scopep{This section defines the evidentiary basis used to formulate the report so that the document
remains auditable.}

\approach{The analysis uses repository-internal evidence only. Source inspection focused on runtime
entry points, public interfaces, build scripts, deployment scripts, documentation generators, and
available tests or operational validation scripts. This avoids introducing unsupported assumptions
about hardware behavior or service contracts that are not visible in the repository.}

\realization{The principal evidence sources are the top-level execution modules
\module{orchestrator.py}, \module{campaign\_runner.py}, \module{status.py}, \module{retry\_queue.py},
\module{server\_webrtc.py}, \module{rf/rf.c}, and \module{gps-lte/gps-lte.c}; the build scripts
\module{CMakeLists.txt}, \module{build.sh}, and \module{install.sh}; the documentation scaffolding
\module{document.sh}, \module{docs/python\_modules.rst}, and \module{docs/c\_modules.rst}; and the
validation assets under \module{test/} and \module{examples/}.}

\chapter{System Foundation}

\section{Operational Problem Definition}
\scopep{The node is designed to function as an embedded spectral-monitoring endpoint that can
receive RF acquisition commands, calculate spectral products, optionally demodulate audio, collect
GPS coordinates, maintain network connectivity, and publish telemetry to remote services.}

\approach{From a theoretical perspective, the system belongs to the class of distributed
instrumentation nodes. Such systems must reconcile three competing concerns: deterministic data-path
execution, flexible supervisory logic, and survivability under intermittent connectivity or partial
hardware failure. The adopted architecture therefore separates hard real-time or near-real-time
operations from policy-driven orchestration.}

\realization{The C subsystem implements the time-sensitive path: HackRF control, IQ ingestion,
ring-buffer management, PSD estimation, demodulation, calibration, GPIO selection, and ZMQ reply
generation. The Python subsystem implements policy-heavy tasks: state transitions, polling the API,
campaign scheduling, retry storage, status reporting, and WebRTC process lifecycle management. Shared
state is maintained through \module{/dev/shm/persistent.json}, exposed in Python through
\module{ShmStore} and in the GPS-LTE C service through dedicated SHM helper functions.}

\section{Architectural Partitioning Principle}
\scopep{This section explains why the repository is intentionally partitioned between C and Python
instead of converging on a single implementation language.}

\approach{The partitioning strategy follows a performance-risk alignment principle. C is used where
the software directly faces high-rate hardware streams, strict memory reuse requirements, callback
contexts, and DSP kernels. Python is used where the dominant concerns are asynchronous orchestration,
serialization, network communication, scheduling, and operator-facing maintainability. This is not
merely a convenience choice; it is a methodological decision that minimizes timing jitter in the data
plane while preserving high development velocity in the control plane.}

\realization{The decision is visible in the repository layout and in the runtime call graph.
\module{rf/rf.c}, \module{rf/libs/psd.c}, \module{rf/libs/chan\_filter.c}, and
\module{rf/libs/sdr\_HAL.c} keep the RF path native and memory-conscious. By contrast,
\module{orchestrator.py}, \module{campaign\_runner.py}, and \module{retry\_queue.py} rely on
Python's dataclasses, HTTP libraries, asynchronous flow control, and OS process management to
implement higher-level behavior.}

\chapter{Architectural Decomposition}

\section{RF Acquisition and DSP Subsystem}
\scopep{This subsystem covers radio hardware control, IQ capture, spectrum estimation, demodulation,
audio streaming preparation, calibration, and response serialization. It is the core measurement
engine of the node.}

\approach{The subsystem follows a staged streaming architecture. A high-frequency hardware callback
performs minimal work and forwards raw IQ samples into ring buffers. Downstream processing then
operates outside the callback context, enabling PSD estimation, demodulation, filtering, and JSON
reply construction without risking upstream sample loss. This is a standard embedded DSP pattern for
decoupling acquisition from computation.}

\realization{The principal implementation is \module{rf/rf.c}. The function \module{rx\_callback}
writes data into the primary ring buffer and conditionally mirrors IQ samples into an audio buffer.
Configuration requests are parsed into \module{DesiredCfg\_t} through \module{parse\_config\_rf} in
\module{rf/libs/parser.c}, mapped to hardware and PSD settings via \module{find\_params\_psd} in
\module{rf/libs/psd.c}, and applied to the HackRF through \module{hackrf\_apply\_cfg} in
\module{rf/libs/sdr\_HAL.c}. The RF engine stores reusable workspaces in dynamic buffers to avoid
repeated allocation costs during repeated acquisitions. Results are serialized by
\module{publish\_results}, which emits PSD arrays and mode-specific metrics such as AM depth or FM
excursion.}

\section{Configuration Integrity for RF Requests}
\scopep{This section addresses how the node prevents malformed or physically inconsistent runtime
requests from propagating directly into the RF pipeline.}

\approach{The design applies two validation layers. The first layer is a typed Python-side request
contract that rejects obviously invalid user or server input. The second layer is a C-side parser
that enforces defaults and clamps filter limits to the Nyquist region derived from the requested
center frequency and sample rate. The dual-layer strategy is methodologically sound because the
Python layer protects normal control flow while the C layer remains authoritative for the actual
measurement engine.}

\realization{On the Python side, \module{ServerRealtimeConfig} in \module{utils/request\_util.py}
validates frequency range, sample-rate bounds, antenna port values, PSD method names, demodulation
mode, and non-negative cooldown intervals. On the native side, \module{rf/libs/parser.c} sets
factory-safe defaults, parses the JSON request, and clamps the optional \module{filter} object so
that its start and end frequencies remain inside the available capture bandwidth. This explicit
clamping is one of the most important correctness barriers in the RF path.}

\section{Calibration and Frequency Error Compensation}
\scopep{The calibration path exists to estimate frequency error and propagate the resulting parts-per-
million correction through subsequent acquisitions.}

\approach{Calibration is treated as a closed-loop estimation problem. The system captures a spectrum,
identifies candidate peaks, evaluates FM pilot-tone evidence, and derives a correction factor that
can be persisted and reused. Methodologically, this is preferable to hardcoded oscillator correction
because it adapts to unit-specific hardware drift and environmental variation.}

\realization{The native calibration routine is implemented in \module{calibrate\_hackrf} inside
\module{rf/rf.c}. Successful calibration concludes with persistence of \module{ppm\_error} into the
shared-memory store. During normal acquisition, \module{fetch\_realtime\_config} in
\module{orchestrator.py} retrieves the current correction from \module{ShmStore}, and
\module{hackrf\_apply\_cfg} in \module{rf/libs/sdr\_HAL.c} computes a corrected tuning frequency using
the standard linear PPM formula. This makes the correction operational rather than merely diagnostic.}

\section{GPS and LTE Subsystem}
\scopep{This subsystem manages modem control, PPP connectivity, GPS acquisition, coordinate
normalization, and publication of mobility information to the external API and local shared state.}

\approach{The design treats connectivity and geolocation as a separate hardware service rather than an
auxiliary function inside the RF process. This separation is methodologically justified because LTE
modem control and GPS parsing involve distinct peripherals, blocking I/O patterns, and failure modes
that should not perturb the RF sampling pipeline.}

\realization{The executable entry point is \module{gps-lte/gps-lte.c}. It contains strict parsing for
numeric coordinates, NMEA-to-decimal conversion helpers, haversine-based displacement checks, and
logic for PPP recovery through \module{connection\_LTE}. The subordinate libraries
\module{gps-lte/libs/bacn\_GPS.c}, \module{gps-lte/libs/bacn\_LTE.c}, and
\module{gps-lte/libs/utils.c} encapsulate GPS UART handling, AT-command transport, HTTP posting, and
shared-memory persistence. The module also writes \module{last\_lat}, \module{last\_lng}, and
\module{changed\_gps} so that location state becomes available to the rest of the node without direct
process coupling.}

\section{Hardware GPIO Abstraction}
\scopep{GPIO support provides board-specific control of modem power, reset behavior, status lines, and
antenna routing.}

\approach{The design uses a minimal hardware abstraction layer around \module{libgpiod}. This is a
sound methodological compromise: the code remains close enough to the platform to preserve exact pin
semantics, but it avoids scattering raw GPIO requests across unrelated application modules.}

\realization{The abstraction is implemented in \module{common/bacn\_gpio.c} and
\module{common/bacn\_gpio.h}. These files centralize pin offsets, line request helpers, modem power
control routines, reset sequencing, status inspection, and antenna selection functions. The build
system can exclude this dependency through \module{BUILD\_STANDALONE}, which preserves local
development on non-target hosts.}

\section{Orchestration and State Management}
\scopep{The orchestration layer determines what the node should do at any moment: remain idle, serve
realtime acquisitions, execute campaigns, or perform calibration.}

\approach{The repository adopts a supervisory finite-state approach. Mutually exclusive high-level
modes are represented explicitly, which prevents accidental overlap between incompatible hardware
operations such as calibration and campaign-driven acquisition. This approach is more robust than an
implicit state model spread across timers and boolean flags.}

\realization{The state machine is centered on \module{SysState} and \module{GlobalSys} in
\module{functions.py}. \module{orchestrator.py} implements the control loop, with
\module{run\_realtime\_logic} handling sticky realtime acquisition and
\module{run\_campaigns\_logic} handling campaign windows. It also manages subordinate processes,
including WebRTC publication, through \module{ManagedProc}, stale-process cleanup helpers, and
explicit stop semantics. The implementation therefore combines logical state management with concrete
process-lifecycle discipline.}

\section{Campaign Scheduling and Deferred Execution}
\scopep{Campaign scheduling is the mechanism by which time-bounded acquisition plans received from the
API are translated into executable local behavior.}

\approach{The design externalizes periodic execution to the operating system's scheduler instead of
implementing its own always-awake timing loop for campaign shots. This is methodologically effective
for embedded Linux systems because it reduces application complexity while leveraging a mature system
scheduler.}

\realization{The scheduler is implemented by \module{CronSchedulerCampaign} in \module{functions.py}.
It synchronizes API campaign metadata, chooses the highest-priority valid campaign, mirrors its key RF
parameters into \module{ShmStore}, and writes a crontab entry pointing to \module{campaign\_runner.py}.
The runner then performs a single acquisition cycle, uploads the result, and applies local storage
policies when network delivery fails.}

\section{Inter-process Communication and Shared State}
\scopep{This section covers the mechanisms by which the separate executables cooperate without being
collapsed into a monolithic process.}

\approach{Two complementary integration mechanisms are used. Request-response control and bulk spectral
results use ZeroMQ over an IPC endpoint, which is well suited to synchronous command execution between
the Python control plane and the C RF engine. Low-frequency persistent node state uses a JSON file in
\module{/dev/shm}, which offers low-latency cross-process visibility while avoiding unnecessary writes
to flash storage.}

\realization{The ZMQ path is implemented by \module{ZmqPairController} in
\module{utils/request\_util.py} and by \module{rf/libs/zmq\_util.[ch]} in the native layer. The
controller enforces strict request-reply semantics and recreates the socket after timeouts to discard
late replies. Shared state is implemented by \module{ShmStore} in \module{utils/io\_util.py}, which
uses file locking and atomic update patterns. The GPS-LTE service uses analogous C helpers in
\module{gps-lte/libs/utils.c} so that both language domains converge on the same persistence model.}

\section{Data Integrity, Retry, and Offline Survivability}
\scopep{This section examines the mechanisms that preserve measurement artifacts during transient API or
network failures.}

\approach{The design follows an eventually-deliverable telemetry model. If a measurement cannot be
uploaded immediately, it should either remain locally recoverable or be discarded only under explicit
resource constraints. The system therefore combines atomic file persistence, bounded queueing, and a
dedicated retry worker.}

\realization{Atomic persistence is implemented by \module{atomic\_write\_bytes} in
\module{utils/io\_util.py}. Campaign samples are stored to \module{Queue/} when uploads fail and to
\module{Historic/} when uploads succeed and disk conditions permit. The dedicated script
\module{retry\_queue.py} reads queued JSON files oldest-first, distinguishes transient and permanent
HTTP failures, deletes corrupt payloads, and preserves ordering by stopping on the first unsent file.
This is a clear and deliberate resilience policy rather than incidental error handling.}

\section{Observability and Operational Telemetry}
\scopep{The node requires a self-reporting capability so that remote operators can assess not only
measurement outputs but also sensor health.}

\approach{Observability is treated as a first-class operational requirement. The node therefore reports
both static identity metadata and dynamic health metrics such as CPU load, memory usage, disk use,
temperature, ping latency, calibration timestamp, and last NTP synchronization. This is methodologically
important because unattended field nodes cannot be managed exclusively through payload success rates.}

\realization{The reporting path is implemented by \module{status.py} and \module{utils/status\_util.py}.
\module{StatusDevice} reads Linux virtual files such as \module{/proc/stat} and \module{/proc/meminfo},
extracts recent log fragments, and assembles a normalized status structure via \module{StatusPost}.
The final payload is transmitted through \module{RequestClient}. Logging itself is hardened by
\module{cfg.py}, which defines \module{AtomicRotator} and environment-sensitive logging filters.}

\section{Audio Streaming and Remote Listening}
\scopep{The system includes an optional audio-monitoring path for demodulated signals, distinct from
the spectral data path.}

\approach{The design uses a split streaming path: native code performs demodulation and Opus-oriented
audio preparation, while a dedicated Python process bridges the encoded audio into a WebRTC session.
This separation is methodologically strong because WebRTC negotiation and GStreamer lifecycle handling
are complex control-plane concerns that should not burden the RF executable.}

\realization{In the native layer, \module{audio\_thread\_fn} in \module{rf/rf.c} drains IQ samples
from a dedicated audio ring buffer, optionally filters them, demodulates AM or FM, and hands frames to
the TCP/Opus transmission support libraries \module{rf/libs/opus\_tx.h} and
\module{rf/libs/net\_audio\_retry.h}. In Python, \module{server\_webrtc.py} accepts the TCP stream,
builds a GStreamer pipeline around \module{webrtcbin}, and maintains signaling with the remote server.
The orchestrator starts and stops this process according to demodulation mode.}

\chapter{Design Rationale and Quality Attributes}

\section{Determinism and Timing Discipline}
\scopep{This section evaluates how the repository preserves timing predictability in the portions of
the node that interact directly with RF hardware and streaming buffers.}

\approach{Determinism is promoted by reducing callback work, using preallocated or reusable buffers,
separating audio and PSD paths, and keeping high-level orchestration out of the RF process. Timing
discipline in SDR systems is less about hard real-time certification than about avoiding systematic
sample loss, reconfiguration races, and unbounded latency spikes.}

\realization{Evidence includes the minimal \module{rx\_callback}, ring-buffer usage in
\module{rf/libs/ring\_buffer.c}, RF processing workspaces in \module{rf/rf.c}, and the request
cooldown mechanism propagated from \module{ServerRealtimeConfig} through \module{parse\_config\_rf}
into the RF engine's runtime state.}

\section{Fault Containment and Recoverability}
\scopep{The system must continue operating or degrade safely when hardware, network, or subordinate
processes fail.}

\approach{Fault containment is achieved through process separation, explicit health checks, bounded
retries, and staged recovery. The methodological aim is not absolute fault elimination, but controlled
failure domains and recoverable degradation.}

\realization{The RF engine includes health checks such as \module{ensure\_hackrf\_session\_is\_healthy}
and \module{recover\_hackrf}. The orchestrator uses stale-process cleanup and managed subprocess
wrappers. Upload failures are handled by queue persistence and retry workers. The GPS-LTE service
contains PPP reconnection logic. These mechanisms collectively show that operational failure was
anticipated as a normal design condition.}

\section{Maintainability and Extension Potential}
\scopep{This section addresses how easily the codebase can be modified for new campaigns, additional RF
modes, or deployment changes.}

\approach{Maintainability in embedded monitoring systems depends on module boundaries more than on
language uniformity. The repository favors explicit modules, small service entry points, configuration
objects, and domain-specific helper libraries. This supports local reasoning and selective replacement.}

\realization{Examples include the typed Python configuration dataclasses, the dedicated DSP libraries
under \module{rf/libs/}, the explicit GPS/LTE helper libraries, and the build-time separation between
standard and standalone modes in \module{CMakeLists.txt} and \module{build.sh}. A future maintainer can
replace a telemetry endpoint, add a new demodulation branch, or alter campaign scheduling with limited
impact on unrelated layers.}

\section{Resource Stewardship in an Embedded Context}
\scopep{The node is intended for resource-constrained deployment, so persistence, logging, and storage
policies must avoid unnecessary wear or uncontrolled growth.}

\approach{The repository consistently adopts flash-aware operational patterns: RAM-backed shared state,
atomic file replacement, rotating logs, bounded retry queues, and selective historical retention. These
choices align with embedded reliability practice, where write amplification and partial-file corruption
are recurring failure sources.}

\realization{The strongest evidence is the use of \module{/dev/shm/persistent.json}, the
\module{AtomicRotator} log manager in \module{cfg.py}, the queue and historic directories managed by
\module{campaign\_runner.py}, and the immediate \module{fsync}-based writes in
\module{utils/io\_util.py}.}

\chapter{Build, Deployment, and Documentation Method}

\section{Compilation Strategy}
\scopep{This section documents how the node is built and how the build structure reflects architectural
intent.}

\approach{The build method distinguishes target deployment from host-side development. This is a
critical methodological feature because hardware-bound projects often stall if local development
requires the full peripheral stack.}

\realization{The top-level \module{CMakeLists.txt} defines a standard build, which produces both
\module{rf\_app} and \module{ltegps\_app}, and a standalone build activated with
\module{BUILD\_STANDALONE}, which removes the GPIO dependency and compiles only the RF application.
\module{build.sh} exposes this distinction through its default mode and its \module{-dev} mode.}

\section{Deployment and Service Registration}
\scopep{The deployment method must provision dependencies, artifacts, shared state, and operating
system services in a reproducible way.}

\approach{The repository uses a scripted installation workflow that approximates infrastructure as
code. Although implemented as shell rather than a configuration-management framework, the script is
methodologically explicit about service shutdown, dependency installation, compilation, permission
repair, shared-memory initialization, and systemd enablement.}

\realization{These responsibilities are captured in \module{install.sh}. The script installs system
packages, optionally compiles \module{libgpiod} and \module{kalibrate-hackrf}, creates the Python
virtual environment, compiles native binaries, initializes \module{/dev/shm/persistent.json},
registers systemd services through \module{init\_sys.py} and \module{daemons/*.service}, and reboots
the system at the end.}

\section{Documentation Pipeline}
\scopep{A mature engineering repository should document not only runtime behavior but also how
documentation artifacts are generated and refreshed.}

\approach{The repository uses a mixed documentation method aligned with its mixed-language
implementation: Doxygen for C-level structural documentation and Sphinx for higher-level rendered
documentation. This is methodologically consistent with the source organization and avoids forcing one
toolchain to understand both language domains equally well.}

\realization{The automation scripts \module{document.sh} and \module{build\_docs.sh} create the
documentation environment, run Doxygen, and build Sphinx outputs. The files
\module{docs/python\_modules.rst} and \module{docs/c\_modules.rst} already define the intended
documentation entry points for the Python and C layers respectively.}

\chapter{Verification Status and Methodological Gaps}

\section{Current Verification Evidence}
\scopep{This section documents what forms of validation are currently present in the repository.}

\approach{Verification evidence is evaluated conservatively. The existence of scripts or examples is
not treated as equivalent to a comprehensive automated test suite; rather, each artifact is classified
according to the level of assurance it can reasonably provide.}

\realization{The repository contains build verification assets, documentation builds, benchmarking
tools, and operational test scripts. Examples include \module{test/test\_benchmarking.py},
\module{test/tester.py}, \module{test/hackrf-max-5min.py}, and the benchmarking support under
\module{utils/benchmarking.py} and \module{examples/benchmark\_rf\_example.py}. These artifacts provide
useful performance and manual-operational evidence, but they do not yet constitute broad unit or
integration-test coverage across the full control plane and native DSP path.}

\section{Limitations Identified During Inspection}
\scopep{A methodological report should record not only strengths but also explicit limitations that
affect future assurance work.}

\approach{The principal limitation is that verification is weighted toward operational scripts and
benchmarking rather than systematic automated assertions over module contracts, failure injection, and
cross-process integration scenarios. This does not invalidate the architecture, but it constrains the
level of confidence that can be claimed from repository-local evidence alone.}

\realization{The most visible gaps are the absence of a broad automated test harness for the RF request
contract, shared-memory concurrency semantics, and orchestrator state transitions under exceptional
conditions. Additional value would come from deterministic tests around \module{ServerRealtimeConfig},
\module{ShmStore}, retry semantics in \module{retry\_queue.py}, and parser boundary handling in
\module{rf/libs/parser.c}.}

\section{Recommended Verification Priorities}
\scopep{This subsection records technically grounded next steps that follow directly from the inspected
architecture.}

\approach{Recommended actions are prioritized by architectural leverage: protect interfaces first,
then validate fault paths, then measure performance regressions. Such sequencing yields faster quality
improvement than beginning with broad end-to-end hardware tests.}

\realization{
\begin{enumerate}
\item add pure-software tests for configuration validation and Nyquist clamping;
\item add concurrency-oriented tests for \module{ShmStore} update semantics;
\item add orchestrator tests that simulate API responses and verify correct state transitions;
\item add contract tests for retry queue behavior under 2xx, 4xx, 5xx, and timeout conditions;
\item preserve the existing benchmarking scripts as regression monitors for throughput and thermal
behavior on target hardware.
\end{enumerate}}

\chapter{Traceability Matrix}

\section{Requirement-to-Module Allocation}
\scopep{The following matrix maps architectural responsibilities to the modules that currently realize
them. It is intended to support maintenance, review, and change-impact analysis.}

\approach{Traceability is organized by engineering concern rather than by folder alone. This improves
its usefulness during audits because reviewers typically reason from required behavior toward source
artifacts, not in the reverse direction.}

\realization{
\begin{longtable}{p{0.22\textwidth}p{0.28\textwidth}p{0.42\textwidth}}
\toprule
\textbf{Engineering concern} & \textbf{Primary modules} & \textbf{Methodological contribution} \\
\midrule
\endhead
RF acquisition and PSD &
\module{rf/rf.c}, \module{rf/libs/psd.c}, \module{rf/libs/ring\_buffer.c} &
Implements the deterministic data path, workspace reuse, ring-buffer decoupling, and PSD generation. \\

Runtime request validation &
\module{utils/request\_util.py}, \module{rf/libs/parser.c} &
Applies typed validation, defaults, and Nyquist-safe clamping before hardware reconfiguration. \\

Frequency calibration &
\module{orchestrator.py}, \module{rf/rf.c}, \module{rf/libs/sdr\_HAL.c} &
Estimates and persists \module{ppm\_error}, then applies corrected tuning frequencies during operation. \\

GPS and LTE telemetry &
\module{gps-lte/gps-lte.c}, \module{gps-lte/libs/bacn\_GPS.c}, \module{gps-lte/libs/bacn\_LTE.c} &
Maintains geolocation and network connectivity in an isolated service domain. \\

GPIO and antenna routing &
\module{common/bacn\_gpio.c}, \module{common/bacn\_gpio.h} &
Centralizes board-level control signals and preserves portability through build-time exclusion. \\

Orchestration and system modes &
\module{orchestrator.py}, \module{functions.py} &
Defines the global state machine, realtime loop, calibration trigger path, and campaign supervision. \\

Campaign execution &
\module{functions.py}, \module{campaign\_runner.py} &
Translates server campaign windows into scheduled local acquisitions and controlled uploads. \\

IPC and shared persistence &
\module{utils/request\_util.py}, \module{rf/libs/zmq\_util.h}, \module{utils/io\_util.py} &
Provides synchronous command exchange and low-wear cross-process state sharing. \\

Offline survivability &
\module{campaign\_runner.py}, \module{retry\_queue.py}, \module{utils/io\_util.py} &
Ensures measurements remain recoverable across network outages and partial API failures. \\

Status and observability &
\module{status.py}, \module{utils/status\_util.py}, \module{cfg.py} &
Collects health metrics, log context, timing metadata, and environment-sensitive diagnostics. \\

Audio publication &
\module{rf/rf.c}, \module{rf/libs/opus\_tx.h}, \module{server\_webrtc.py} &
Supports remote listening while isolating signaling complexity from the RF executable. \\

Build and deployment &
\module{CMakeLists.txt}, \module{build.sh}, \module{install.sh} &
Encodes reproducible compilation and target provisioning with explicit dev/target separation. \\
\bottomrule
\end{longtable}}

\chapter{Conclusion}

\section{Synthesis}
\scopep{This concluding section consolidates the methodological interpretation of the repository as a
whole.}

\approach{The node should be understood as a deliberately layered sensing platform rather than as a
single application. Its architecture reflects an embedded-systems design method in which hardware-
proximate, timing-sensitive, and buffer-heavy responsibilities remain in C, while coordination,
policy, networking, and operator-facing workflows remain in Python. The integration strategy relies on
simple, inspectable mechanisms: synchronous IPC over ZeroMQ and RAM-backed JSON persistence.}

\realization{The repository expresses this method consistently across its implementation and tooling:
the RF engine, GPS-LTE service, orchestration loop, campaign runner, retry worker, status reporter,
and documentation pipeline are distinct but interoperable artifacts. The strongest architectural
qualities evidenced by the code are separation of concerns, resilience to intermittent failures,
flash-aware persistence discipline, and a clear path for incremental enhancement. The main outstanding
engineering need is broader automated verification around contract boundaries and failure scenarios.}

\appendix

\chapter{Key Source Inventory}

\begin{tabularx}{\textwidth}{p{0.28\textwidth}Y}
\toprule
\textbf{Path} & \textbf{Primary role} \\
\midrule
\module{rf/rf.c} & Main RF executable, acquisition loop, calibration, audio thread, JSON reply generation. \\
\module{rf/libs/parser.c} & JSON parsing, defaults, and boundary enforcement for RF requests. \\
\module{rf/libs/psd.c} & Spectral estimation algorithms and parameter derivation. \\
\module{rf/libs/sdr\_HAL.c} & HackRF hardware configuration with PPM-aware tuning. \\
\module{gps-lte/gps-lte.c} & GPS acquisition, LTE connectivity management, and coordinate reporting. \\
\module{common/bacn\_gpio.c} & GPIO control for modem power/reset and antenna selection. \\
\module{orchestrator.py} & Global supervisor for realtime and campaign modes. \\
\module{functions.py} & Shared orchestration helpers, state model, scheduler, acquisition wrapper. \\
\module{campaign\_runner.py} & One-shot campaign acquisition and upload worker. \\
\module{retry\_queue.py} & Deferred upload recovery worker. \\
\module{status.py} & Status payload generation and publication entry point. \\
\module{server\_webrtc.py} & WebRTC/GStreamer bridge for remote audio monitoring. \\
\module{utils/io\_util.py} & Atomic writes, shared-memory persistence, timers. \\
\module{utils/request\_util.py} & HTTP client, ZMQ controller, configuration dataclasses. \\
\module{utils/status\_util.py} & Host-health telemetry collection and shaping. \\
\module{cfg.py} & Environment loading, paths, identity helpers, and logging infrastructure. \\
\module{CMakeLists.txt}, \module{build.sh}, \module{install.sh} & Build and deployment baseline. \\
\bottomrule
\end{tabularx}

\end{document}
